\documentclass[11pt,a4paper]{article}

\usepackage{fontspec}
\setmainfont{Arial}

\usepackage{geometry}
\geometry{a4paper,top=2cm,bottom=2cm,left=2cm,right=2cm}

\usepackage[ngerman]{babel}

\usepackage{titlesec}
\titleformat{\section}{\fontsize{12}{14.4}\selectfont\bfseries}{\thesection}{1em}{}
\titlespacing*{\section}{0pt}{8pt}{4pt}

\setlength{\parindent}{0pt}
\setlength{\parskip}{6pt}

\usepackage{fancyvrb}
\fvset{fontsize=\small,frame=single,framesep=3pt,xleftmargin=2pt,xrightmargin=2pt}

\usepackage[hidelinks]{hyperref}

\begin{document}
\pagestyle{empty}

{\LARGE\bfseries Installationsanleitung: Uni Speedrun}

\textbf{Repository:} \url{https://github.com/FreakyLetsFail/UNI-Speedrun}

\section*{Voraussetzungen}
Windows 10/11, Python 3.10 oder neuer, Git und Windows Terminal oder Visual Studio Code.

\section*{Installation unter Windows}
\begin{Verbatim}
git clone https://github.com/FreakyLetsFail/UNI-Speedrun.git
cd UNI-Speedrun
py -m venv .venv
.\.venv\Scripts\Activate.ps1
python -m pip install --upgrade pip
python -m pip install -r requirements.txt
\end{Verbatim}

\section*{Start}
\begin{Verbatim}
$env:PYTHONPATH = "src"
python -m uni_speedrun
\end{Verbatim}
Beim ersten Start wird die lokale Datenbank unter \texttt{data/uni\_speedrun.db} angelegt. Anschließend kann ein Studienplan erstellt werden.

\section*{Tests}
\begin{Verbatim}
python -m pip install -r requirements-dev.txt
python -m pytest
\end{Verbatim}
Erwartetes Ergebnis: 56 erfolgreiche Tests.

\end{document}
